El acuerdo global es de \textbf{4.47 puntos porcentuales} de RMS sobre los ocho puntos. De los 3 que tienen réplicas, 0 caen dentro de dos desviaciones de la medida.

\textbf{La meseta.} El modelo se estabiliza en \qty{28.70}{\percent} y el ensayo en \qty{25.20}{\percent}. Descontada la humedad, eso significa que el modelo libera el \qty{106}{\percent} de la materia volátil del análisis próximo y el ensayo el \qty{93}{\percent}. No es un detalle numérico: dice que en el ensayo parte del volátil \emph{no llega a salir}, algo esperable porque el análisis próximo se hace a \SI{950}{\celsius} durante \SI{7}{\minute} en crisol tapado, y el ensayo es otra cosa.

\textbf{Lo que el contraste NO puede separar.} Un residuo en un instante dado admite dos lecturas: que la carga esté a otra temperatura o que la cinética sea otra. Con la curva de masa sola no se distinguen, porque sólo se observa el producto de ambas. La cronología del aglomerado —hinchamiento a los \SI{90}{\second}— sí las separa, porque el hinchamiento depende de la \emph{temperatura} y no de la cinética de liberación: fija la historia térmica y deja la cinética como única incógnita. Ese es el argumento que llevó a la compuerta y no al acoplamiento térmico.
